% ECONOMICS_PROBLEM: {"id": "C78-268", "title": "House exchanges along longer social-network cycles", "jel_code": "C78", "source_id": "OP-0550"}
% FIRST_POSTED: 2026-10-10
% ATTEMPT_STATUS: solved
% ENTRY_KIND: research_attempt
\documentclass[11pt]{article}
\usepackage[a4paper,margin=25mm]{geometry}
\usepackage[T1]{fontenc}
\usepackage{lmodern,microtype,amsmath,amssymb,amsthm}
\usepackage[hidelinks]{hyperref}
\newtheorem{theorem}{Theorem}
\newtheorem{lemma}{Lemma}
\newcommand{\NP}{\mathsf{NP}}
\newcommand{\FP}{\mathsf{FP}}
\newcommand{\RA}{\mathcal R}
\title{Reachability and Pareto selection with longer network trades}
\author{Ji Ho Bae}
\date{10 October 2026}
\hypersetup{pdftitle={Reachability and Pareto selection with longer network trades},pdfauthor={Ji Ho Bae}}
\begin{document}
\maketitle

\paragraph{Authorship and provenance.}
Author: Ji Ho Bae. Prepared with GPT Codex assistance; the precise serving
revision was not exposed. The complete original argument is retained in
the \href{https://github.com/jbaelaw/long-cycle-trade-reachability/blob/6ff34b79452d2e25ee6490852608610a949f911b/paper/Proof.tex}{immutable TeX source} and
\href{https://github.com/jbaelaw/long-cycle-trade-reachability/blob/6ff34b79452d2e25ee6490852608610a949f911b/paper/Proof.pdf}{hand-proof PDF}. This submission adds catalogue metadata,
provenance and a completion statement to that argument. Its preparation
included exact-target and primary-source checks and separate mathematical
reviews. This is a full hand-proof claim submitted for independent expert
review; community acceptance remains pending.

\begin{abstract}
We classify improving cyclic trades on an undirected social network.
For every fixed length bound $k\geq3$, for an input bound $k\geq3$,
and for unrestricted length, both allocation reachability and agent--house
reachability are NP-complete. Selecting an allocation that is Pareto
optimal within the entire reachable set is polynomial-time Turing
equivalent to an NP-complete decision oracle. A private-house invariant
prevents longer trades from invalidating the completion reduction.
\end{abstract}

\section{Model and result}
There are $n$ agents, $n$ houses, an undirected graph on the agents, and
an initial bijection. Write $M(i)$ for agent $i$'s house in an allocation $M$.
Each agent strictly orders her acceptable houses,
which include her endowment. A trade uses distinct agents
$a_1,\ldots,a_\ell$, with $2\leq\ell\leq k$, consecutive graph edges
including the closing edge, and gives $a_i$ the current house of
$a_{i+1}$, cyclically. Every participant must strictly improve.
A two-agent trade uses one graph edge. Unrestricted length means $k=n$;
an input bound is effectively capped at $n$.

Let $\RA$ be all endpoints reachable from the initial allocation.
A member of $\RA$ is Pareto optimal in $\RA$ if no other member weakly
improves every agent and strictly improves at least one. The comparison
is with all of $\RA$, not just the continuations of the selected endpoint.
These are the network-edge conventions in catalogue problem C78-268.
The motivating survey proposes longer-cycle exchanges as a research
direction without specifying these conventions~\cite[\S3.1.3]{survey}.

\begin{theorem}\label{main}
The following statements hold for each fixed $k\geq3$, for an input
bound $k\geq3$, and for unrestricted length.
\begin{enumerate}
\item Reachability of a specified allocation is NP-complete.
\item Reachability of a specified house for a specified agent is NP-complete,
      even on trees.
\item Reachable Pareto selection admits an $\FP^{\NP}$ algorithm.
      Conversely, every $\FP^{\NP}$ function can be computed in polynomial
      time with an oracle returning any valid reachable Pareto allocation.
\end{enumerate}
The decision reductions are polynomial-time many-one reductions. The
search equivalence uses polynomial-time Turing reductions that succeed
for every valid oracle answer. In particular, each of the three tasks
has a polynomial-time algorithm if and only if $\mathsf P=\NP$.
\end{theorem}

\section{Short paths and upper bounds}
Give agent $i$ integer ranks $r_i(h)\in\{0,\ldots,n-1\}$ increasing
with preference. Every trade increases $\sum_i r_i(M(i))$ by at least
two, while this sum is at most $n(n-1)$. Every legal sequence therefore
has at most $\lfloor n(n-1)/2\rfloor$ trades. A certificate lists these
cycles; distinctness, graph edges, acceptability, improvement and the
endpoint can all be checked in polynomial time. This proves both NP
upper bounds in all three regimes.

For Pareto selection, fix agents in order. For the current agent,
test her houses from best to worst, asking whether a reachable endpoint
assigns that house and all previously fixed houses. Each question is
in NP by the path bound. At least one answer is positive, since the
previous choices preserved a reachable endpoint. After at most $n^2$
queries, all agents' houses are fixed. If a reachable allocation Pareto
dominated the result, its first different agent would have received a
better feasible choice at her selection step. This is a contradiction.
All queries start from the original endowment. A witnessing path can
also be recovered by polynomially many certificate-prefix queries.

\section{A completion gadget that blocks mixed cycles}
We use the following established hardness result: under improving
\emph{swaps}, agent--house reachability is NP-complete on trees
\cite[Theorem 1]{swaps}. Those hard instances have a target house
initially held by a different agent. Relabel the target agent as $1$
and write the original house set as $X=\{x_1,\ldots,x_n\}$ and the target as
$t\ne x_1$. A tree has no simple cycle of length at least three,
so this theorem already proves part~2 of Theorem~\ref{main}.

To reduce this tree problem to allocation reachability, keep its agents
and edges. For each original agent $j$, add a copy $c_j$ endowed with a
new house $d_j$, and add the edge $\{j,c_j\}$. Make the copies a clique;
there are no other edges between the two parts. Original agent $j$
ranks $d_j$ first, then all original houses in her original order,
then all other private houses. These lists may be completed strictly;
all original instances used here already have complete lists.

Copy $c_1$ ranks $t$ first, then $d_1$, then all other houses. List
$X\setminus\{t\}$ as distinct houses $t_2,\ldots,t_n$ in any fixed
order. For $j=2,\ldots,n$, copy $c_j$ ranks $t_j$ first, then the other
houses in the list $t_n,t_{n-1},\ldots,t_2$, then $d_j$, then $t$ and
all foreign private houses. Order any otherwise unspecified houses arbitrarily.
The distinct top choices define one allocation
\[
 T(j)=d_j\quad(1\leq j\leq n),\qquad
 T(c_1)=t,\qquad T(c_j)=t_j\quad(2\leq j\leq n).
\]
The construction has $2n$ agents and polynomial encoding size. It
modifies the copy-clique completion used for swaps in
\cite[Theorem 2]{swaps}; the next lemma is needed because longer trades
are now allowed.

\begin{lemma}\label{cut}
Every legal trade of length at least three in the constructed instance
uses only copies. A trade between the two parts is necessarily the swap
of $j$ with $c_j$, after which $j$ never trades again.
\end{lemma}
\begin{proof}
Every agent other than $j$ and $c_j$ ranks $d_j$ below her endowment. Strict improvement keeps every agent's holding
at least as good as her endowment. Thus $d_j$ can leave $c_j$ only for
$j$. On receiving $d_j$, agent $j$ obtains her top choice and never
trades again. In particular, whenever $j$ participates in a trade,
$c_j$ still holds $d_j$.

Consider a legal cycle of length at least three that meets both parts.
It must transfer a house from an original agent $j$ to a copy, which
can only be $c_j$. This copy still holds $d_j$ and must give it to
its other neighbor in the cycle, a different copy: returning it to $j$
would repeat a vertex and make the cycle a two-agent swap. That recipient
cannot improve by receiving $d_j$, a contradiction. No mixed cycle is
legal. The original part is a tree, so it contains no remaining cycle
of length at least three. Finally, a crossing swap gives $j$ the house
$d_j$ and hence retires $j$ permanently.
\end{proof}

\begin{lemma}\label{equiv}
The original tree instance can deliver $t$ to agent $1$ if and only if
$T$ is reachable in the constructed instance, for every allowed length
bound at least two.
\end{lemma}
\begin{proof}
Suppose first that original swaps deliver $t$ to $1$. Perform the same
swaps in the original part, then swap each $j$ with $c_j$. Original
agents obtain their top choices. Copy $c_1$ obtains $t$; every other
copy receives an original house different from $t$ and therefore improves
from her private house.

The remaining copies hold a permutation of $t_2,\ldots,t_n$.
Process $c_2,\ldots,c_n$ in this order. If $c_j$ does not hold $t_j$,
swap it with its holder $c_s$. Previously processed copies already hold
their distinct targets, so $s>j$. The house currently at $c_j$ is some
$t_q$ with $q>j$. Copy $c_j$ receives its top choice. Copy $c_s$ prefers
$t_q$ to $t_j$: either $q=s$ gives its own top choice, or both houses
are ordered by the common descending index order. The swap is legal.
At the end the allocation is $T$. Only swaps were used.

Conversely, suppose a sequence reaches $T$. A copy other than $c_1$
never holds $t$, since she ranks it below her initial house. The first
transfer of $t$ into the copy part must therefore deliver it to $c_1$.
By Lemma~\ref{cut}, this is the swap of $1$ with $c_1$, so agent $1$
held $t$ just before that swap.

Retain only the original--original swaps preceding this event.
They form a legal sequence in the original tree. To see this inductively,
ignore each copy-only trade and each crossing swap. A crossing swap
retires its original participant, who never appears in a later retained
swap; in the projected sequence she simply keeps the original house
she surrendered. Every still-active original agent has exactly the same
house in the projected and enlarged sequences. Hence each retained swap
has the same two houses and remains improving. The projected sequence
therefore delivers $t$ to $1$.
\end{proof}

All newly allowed longer cycles are covered by Lemma~\ref{cut}, even
when the bound is input or unrestricted. Lemma~\ref{equiv} and the
source tree hardness prove part~1 of Theorem~\ref{main}.

\section{Pareto search hardness}
Every agent has a distinct top house in the constructed instance.
If the source instance is positive, $T$ is reachable and Pareto dominates
every different allocation, so every valid reachable Pareto answer is
$T$. If the source instance is negative, $T$ is unreachable, so every
valid answer differs from $T$. One oracle call and an equality test
therefore decide the NP-complete tree problem. This is the all-top
argument for swaps in \cite[Proposition 4]{swaps}, with longer-cycle
soundness supplied by Lemma~\ref{cut}.

More precisely, let $f$ be computed by a polynomial-time transducer
with an NP oracle. Reduce each oracle question to the NP-complete tree
problem, construct the gadget, call a reachable Pareto oracle and test
whether its answer is $T$. The answer is correct regardless of which
valid Pareto allocation the oracle returns. Replacing every adaptive NP
query in this way computes $f$ in polynomial time. Together with the
selector already proved, this establishes part~3 in the stated Turing
sense. It does not assert polynomial-time verification of arbitrary
Pareto answers or completeness under a stronger reduction notion.

\section*{Completion Estimate}
100\%
The stated theorem covers allocation reachability, agent--house reachability
and Pareto selection for every length regime in the catalogue question.
This is the author's complete hand-proof claim, pending independent review.

\begin{thebibliography}{9}
\bibitem{swaps}
L. Gourv\`es, J. Lesca and A. Wilczynski.
Object Allocation via Swaps along a Social Network.
\emph{Proceedings of IJCAI 2017}, 213--219.
\href{https://doi.org/10.24963/ijcai.2017/31}{doi:10.24963/ijcai.2017/31}.
\bibitem{survey}
K. Cechl\'arov\'a, \'A. Cseh and D. F. Manlove.
Selected open problems in Matching Under Preferences.
\emph{Bulletin of the EATCS} 128 (2019), 14--38, \S3.1.3.
\href{https://hpi.de/friedrich/docs/publications/2019/BEATCS.pdf}{Primary manuscript}.
\end{thebibliography}
\end{document}
